\subsection{Model Validations}

The reimplementation of the five models has worked to varying degrees. However, an overall trend was noticed during the reimplementation process: the documentation of a model is crucial for the success of the reimplementation. The complexity, or the number of quantities in the model, does not seem to be a factor in the reimplementation's success. This can be seen in the Saadat2021 model, which is by far the largest model~\cleverref{Table}{tab:models-meta}, but, with the Matuszynska2016 model, is the most successful reimplementation. On the other hand, the Li2021 model, which did not include any information in the publication itself, had not only the most difficult reimplementation process but also the most issues in the results. This shows the need for transparent documentation of the model creation process to enable its use beyond publication. The issues with the model reimplementations are discussed in more detail below.

\subsubsection{Bellasio2019}

The biggest issue with the reimplementation of the Bellasio2019 model is the \gls{ca}, proven primarily by the lacking reproduction of the fifth figure of the publication~\figref{fig:bellasio2019-fig6}. This figure shows the model's response to the transition of \gls{ca}. In both the low-to-high and high-to-low transitions, none of the results match those reported in the publication. Additionally, as the results only change slightly after the transition, it is clear that the model is not responding to the \gls{ca} changes as it should. Strange as the quasi-steady-state scans of both the rates and concentrations, match the publication pretty well in the regard of \gls{ca}~(see Fig.~\ref{fig:bellasio2019-fig3} and Fig.~\ref{fig:bellasio2019-fig4}, both on the right). However, because the x-axis values in both figures match those in the publication, there may be an issue with how the sixth figure was generated. During implementation, it was difficult to follow the model documentation, as there were discrepancies between the text and the equations in the provided Microsoft Excel Workbook. Because of this, and the only minor inconsistencies in the other figures, the model's implementation is classified as a partial success. The math behind the model, especially concerning the \gls{ca} dynamics, need to be reevaluted and perhaps corrected to match the implementation used in the figures of the publication.
\subsubsection{Fuente2024}

Overall, the basic mechanisms of the Fuente2024 model are reproduced, as many results match those reported in the publication. However, a distinct issue is evident at higher frequencies, resulting not only in results shifted away from the publication but also in different dynamics. Further, the results of the \gls{F} could not be reproduced at all~\figref{fig:fuente2024-fig5}, which made it impossible even to attempt to reproduce three of the figures in the publication. The \gls{F} of the reimplementation shows little to no reaction to the oscillations of the light. In contrast, the publication shows a clear curve.

Great care was taken to ensure the reimplementation of the model was as close as possible to the publication, drawing heavily on the provided Wolfram Notebook. However, differences in the software used during implementation may have caused the issues. This is only an assumption, as the exact details of the issue could not be determined. Still, the reimplementation is classified as a partial success, as most of the results closely match the publication. The issues with the \gls{F} and the high frequencies need to be investigated further to determine whether they are caused by a software issue or by the model itself.

\subsubsection{Li2021}

Most of the results of the Li2021 model are reproduced well, however, only for the lower light intensity~\figref{fig:li2021-fig4}. The higher the light intensity, the more the results deviate from the publication~\figref[, top row]{fig:li2021-fig5}. The trend of each result is similar, but does not exactly match the publication. Because the default value of the \gls{ppfd} is \qty{100}{\micro\mol\per\meter\squared\per\second}, the issues in the higher light intensity most likely stem from an equation in the model that is based on the \gls{ppfd}. Additionally, no information about the model is given in the publication itself; only changes made to a preexisting model and a link to a GitHub repository that includes a Python script supposedly implementing the model.

This script is not well documented, written, or formatted. It was the only source of information for the reimplementation of the model, making it easy to make mistakes in the code. Additionally, running the script with the default parameters, except for \gls{ppfd} being set to \qty{500}{\micro\mol\per\meter\squared\per\second}, does not produce the results shown in the publication. These results, however, are very similar to those of the reimplementation; therefore, the script provided is either incomplete or not the actual implementation used for the publication. Because of this, the reimplementation of the model is considered a success, as it reproduced most of the publication's results, albeit with some issues. However, as no other information about the model is provided, it is impossible to investigate the issues in any other way. Additionally, some figures in the publication are not well-explained, making it difficult to reproduce them. The authors should therefore be contacted to ask for more information on the model and its implementation.

\subsubsection{Matuszynska2016}

The reimplementation of the Matuszynska2016 model was overall successful, as all but one figure could be reproduced very well. This outlier was a phase-contrast projection of lumen pH versus \gls{Q} at three different light intensities~\figref[, bottom]{fig:matuszynska2016-fig5}. As these inconsistencies are largely in the y-axis direction, it is fair to assume the lumen pH is the issue, which is reinforced, as the time series of the lumen pH also shows some, but small, inconsistencies~\figref[, top]{fig:matuszynska2016-fig5}. The exact cause of these errors remains unknown, despite careful investigation of the model. It has to be added that the lumenal pH is calculated through the lumenal \gls{h} concentration, blowing up a rather small number. Therefore, a small error in concentration can lead to a large error in pH. Still, the reimplementation is considered a success, as the model largely reproduces not only the figures in the publication but also the experimental data the authors used to validate it. The issues with the lumenal pH need to be investigated further, to determine whether they are caused by a software issue or by an issue with the model itself.

\subsubsection{Saadat2021}

Overall, the reproduction of the Saadat2021 publication was successful, with most figures reproduced well. The only issue is with specific ranges of the \gls{kcyc} that prevent the model from reaching a steady state. A feature that is presented in the publication, but could not be reproduced well in the reimplementation~\figref{fig:saadat2021-fig3}. If this issue is addressed, the model would be able to reproduce all the figures in the publication.

\subsection{Model Demonstrations}

\subsubsection{Daylight Simulations}

In the daylight simulations, the models are simulated using actual daylight data to see how they perform in a more realistic "environment". The plots showcase three simulation results: the \gls{vc}, the \gls{atp} to \gls{nadph} ratio, and the \gls{F}. The five models vary greatly in their results, with some showing a clear response to daylight changes, while others exhibit more consistent behavior. Some even have no result at all~\figref{fig:daylight-demon}.

The \gls{vc} is only shown by the Bellasio2019 and Saadat2021 models. Both models show a proportional correlation between the \gls{vc} and the daylight changes, with the Bellasio2019 model showing a more pronounced response. On top of that, they both reach a plateau before the highest light intensity is reached, both at approximately the same time. This means that \gls{vc} is influenced by \gls{ppfd} in both models, but only up to a certain level. Additionally, the representation of \gls{vc} in both models is likely very similar, as they exhibit a similar response.

Only the Saadat2021 model shows a clear response in the \gls{atp} to \gls{nadph} ratio, with a slight negative correlation with daylight changes; however, the plot remains mostly constant. This can be interpreted as meaning that the rates influencing \gls{atp} and \gls{nadph} are either not influenced by changes in light, or they are influenced in a way that cancels each other out.

The \gls{F} results for the Fuente2024 and Matuszynska2016 models show a clear, very similar response to daylight changes, with a positive correlation. Additionally, both models show the curve closely follows daylight changes at higher light intensities, while showing a plateau at lower light intensities. This means that in both models, the \gls{F} has a minimal base level that is rarely crossed, no matter the light intensity, but easily increases when the light intensity is high enough. This interpretation is supported by the fact that the Fuente2024 model calculates the \gls{F} with a base level implemented, and the Matuszynska2016 model uses a function based on two states of \gls{psii}, both of which are dependent on each other. The only other model showing a response in the \gls{F} is the Saadat2021 model, which shows a completely different response, with a negative correlation to the daylight changes. While this reaction is counterintuitive, it actually follows data from diurnal experiments~\cite{oivukkamakiFieldIntegrationShoot2024}. It was shown that \gls{F} actually falls after a certain light intensity, reaches a lower plateau, and rises again once the light intensity is lower. A very similar curve to the one shown by the Saadat2021 model. This could be explained by the fact that the \gls{npq} mechanism is activated at higher light intensities, leading to a decrease in \gls{F}. This would mean that the Saadat2021 model better describes the \gls{npq} phenomenon than the other models, even though all implement a \gls{npq} mechanism.

The Li2021 model shows no response to changes in daylight in any of the three results. This is not because the simulation could not be done, but because that model does not contain any of the quantities plotted. This shows a clear issue with this demonstration, as it does not apply to all models. A possible solution to this issue would be to add additional quantities to be plotted, or at least present a visual result to demonstrate that the model is responding to daylight changes, even if the specific quantities plotted are not present in the model.

\subsubsection{\glsxtrshort{fvcb} Add-on}

Only the Bellasio2019 and Saadat2021 models have quantities representative of \gls{co2} and \gls{vc}, thereby allowing the \gls{fvcb} to be added in~\figref {fig:fvcb-demon}. The Bellasio2019 model shows a very similar trend in both curves to the \gls{fvcb} curves, actually matching at lower \gls{ci} pressures. In this range, the \gls{fvcb} curve with the default parameters has the \gls{wc} as the limiting factor, which means the Bellasio2019 model can capture this limitation pretty well. However, as the \gls{ci} pressure increases, the limiting factor switches to \gls{wj} and then \gls{wp}, which is not shown in the Bellasio2019 curves. Still, the Bellasio2019 model has a good relationship with the \gls{fvcb} model, even if the curves include some jaggedness at some points. These happen because the model has difficulty simulating a steady state. Adding the \gls{fvcb} to the Bellasio2019 model was easy, since it already includes its simplest form of calculating the \gls{A}~\cleverref[, top]{Equation}{eq:fvcb-A}.

The Saadat2021 model has its own interpretation of \gls{vc}, which was taken to calculate the \gls{A} with the default parameters of the \gls{fvcb} model. The resulting curves are close to those of the \gls{fvcb} model but slightly higher. As both the \gls{vc} and the \gls{A} curves are higher, the reason for the shift can be directly attributed to the \gls{vc}. Both the \gls{vc} and the \gls{wc} of the Saadat2021 and \gls{fvcb} models use Michaelis-Menten kinetics, yet the parameters differ. Additionally, the \gls{vc} of the Saadat2021 model is much more complex and includes many more quantities than the \gls{wc} of the \gls{fvcb} model. This means that the \gls{vc} of the Saadat2021 model is vastly different than the \gls{wc} of the \gls{fvcb} model. However, as the relationship of both curves shows similar trends to the \gls{fvcb} curves, it can be said that the addition of the \gls{fvcb} to the Saadat2021 model was successful, even though a more detailed investigation into the equations and the ways to bridge the two models should be done.

The \gls{fvcb} model could not be added to the other three models because they do not include the necessary quantities. The Matuszynska2016 and Fuente2024 models do not include \gls{vc} or \gls{co2} related quantities. In contrast, the Li2021 model includes a very simplified term for the \gls{cbb} cycle but also lacks a quantity related to \gls{co2}. The fact that these quantities are missing makes it hard to incorporate the simple \gls{fvcb} model, yet it does not degrade the models at all. They all have a different scope of photosynthesis, and these three models focus more on reactions not related to the \gls{cbb} cycle.

\subsubsection{Standard \glsxtrshort{pam} Simulation}

Only the Bellasio2019 model could not be used for the standard \gls{pam} simulation~\figref{fig:pam-demon}, as it does not include a quantity for \gls{F} or \gls{npq}. All other models, however, produced very similar outputs. The Fuente2024, Matuszynska2016, and Saadat2021 models all show a very similar reaction for \gls{F}, with the same number of \gls{Fm} and similar responses to light changes. The Matuszynska2016 and Saadat2021 models use these \gls{F} to calculate the \gls{npq}~\cleverref{Equation}{eq:npq}, resulting in jagged curves for both. However, the Fuente2024 model provides a dynamic description of \gls{npq}, allowing it to show a smooth curve that also shows that \gls{npq} reacts to the strong saturating pulse, a feature that cannot be shown in the other two models. While the Li2021 model does not have a quantity for \gls{F}, it does have a term for the \gls{npq}, which also showcases a similar curve as the other three models, with the added benefit of showcasing the reaction of \gls{npq} to the saturating pulses.

All possible simulations of the standard \gls{pam} protocol show a commonly observed trend in experiments~\cite{matuszynskaMathematicalModelNonphotochemical2016,niesChlorophyllFluorescenceHow2021}, even the Li2021 model, which had difficulty simulating the \gls{npq}, as shown in the publication. These new results of the \gls{npq}, match up much more than the \gls{npq} seen at a light intensity of \qty{500}{\micro\mol\per\meter\squared\per\second}~\figref{fig:li2021-fig3}, which can be interpreted that the red parts of the figures in the publication are actually run at a light intensity of \qty{1000}{\micro\mol\per\meter\squared\per\second}. However, to be certain, a deeper investigation into the model and the publication is needed.

The results of this demonstration, as with the others, cannot be quantified or directly compared. As all models are parametrized differently and have different scopes, it is not expected that their results will match. However, as the \gls{F} of the Fuente2024, Matuszynska2016, and Saadat2021 models are very similar, it can be said that these three models have a similar description of this quantity.

\subsubsection{\glsxtrshort{mca} of Photosynthesis}

Many of the models show an incomplete heatmap for the \gls{mca} of photosynthesis~\figref{fig:mca-demon}, as they do not include all the quantities chosen to be integral for the depiction of photosynthesis. In fact, even the model showing the most responses, the Saadat2021 model, is missing two quantities. The \gls{co2} and a parameter representing the coefficient of \gls{psi}. The former is considered lacking in this regard, as it is included in the model as a parameter rather than a variable. The latter could not be assigned a value because no definitive proportional parameter was found for the rate of \gls{psi} in the model. The heatmaps presented by this model show little response. On the variables side, only the parameter representing \gls{rubisco} shows a control on \gls{rubp}.

In contrast, all the rest show little control over the variables. On the fluxes side, the parameter for \gls{psii} shows some control on all the fluxes. This suggests that the chosen representatives of the photosynthetic machinery do not have much control over one another, except for \gls{psii}.

The other models have very little representation in the heatmaps because they do not include many of the chosen quantities. Therefore, commenting on the system's control is difficult. On top of that, the Li2021 model shows no response in the heatmaps; however, since some labels are not grayed out, it is not because there were no quantities. As an \gls{mca} of control coefficients is done by simulating to steady-state, these heatmaps directly show that the Li2021 model does not simulate to steady-state.

While this \gls{mca} demonstration is not a good way to assess the control of the system, it does directly enable comparison of the availability of the photosynthesis machinery across the different models. With the grayed-out labels and empty squares, the user can, with one glance, see which parts are included in the model and which are not. In this case, the Saadat2021 model covers most of photosynthesis; therefore, if a user is looking for a more complete model, they choose that one. On the other hand, if a user is looking for a more focused model, for example, for \gls{psii} and \gls{rubisco}, they can choose the Bellasio2019 model. A bonus is that the user can also see which model actually simulates to steady-state, as the Li2021 model does not, which can be an important factor for some users.

\subsubsection{Fitting of \glsxtrshort{npq}}

As the Bellassio2019 model does not have a quantity for \gls{F} nor \gls{npq}, it is the only model that cannot be used for the fitting~\figref{fig:fit-demon}. All other models were fitted to the experimental data; however, to varying degrees of success. The best fit is with the Li2021 model, while the Matuzynska2016 model is also a strong contender. The former, however, does not have a quantity for \gls{F}, while the \gls{F} result of the Matuzynska2016 fit shows a very similar pattern to the data. The worst fit is done by the Saadat2021 model, even though the \gls{F} fits better here than with the Fuente2024 model. However, as all models show a better fit to the data than the normal parameterization, it can be said that all these models can be fitted to that data.

The problem with this demonstration is that it is very difficult to find the best way to fit these models. A fitting routine is strongly dependent on the chosen parameters and whether they have bounds. Additionally, the fitting routine itself can be based on different algorithms, leading to different results. Therefore, this demonstration cannot be used to directly compare the models, as the fitting process cannot be standardized across models. However, it can give an idea of what to do with a model and create the starting blocks for a fitting process, which can be very useful for users who are not familiar with fitting models to data.

\subsection{Website}

As the internet is now an integral part of most societies worldwide, everyone knows how to use a website~\cite{MeasuringDigitalDevelopment}. It keeps it simple and accessible for everyone, as it only requires an active internet connection. On top of that, it is a quick way to share information worldwide, as direct sharing only requires a link. However, it is vital to give a website a good design tailored to the target audience. This ensures the best way to convey the information it intends to convey~\cite {DataReportalGlobalDigital}. The advantages of a website are clear: accessibility, ease of use, and quick sharing. Advantages that support the goal of \greensloth{}.

The website created for \greensloth{} is designed to be as simple as possible, with a clear structure and easy navigation~\figref{fig:website-landingpage}. It follows a modern, common style to convey its message while separating key aspects of the work into distinct parts using models. The lists of implemented models and the ability to search for models by custom tags~\figref{fig:website-modelpage} can help the user narrow their focus to the model they wish to look at. This would alleviate the need for an extensive literature search to find a model. After selecting a model, the user can look at all the details of the model, all on one page~\figref {fig:website-model-top}. This includes the model equations, parameters, results, and a link to the code. This allows the user to quickly get an overview of the model and decide whether it is relevant to their work. As the information is directly sourced from the GitHub repository documentation, it is always up to date with the latest code changes. Creating a direct bridge between the code and the website.

From each model page, the user can also navigate to the comparison page, where they can compare all the models on \greensloth{}~\figref{fig:website-compare-var}. There, they choose the models and the aspects they want to compare. By choosing between Variables, Simulation Examples, Model Information, Demonstrations, and Schemes, they can get a good overview of the differences and similarities between the models. These comparisons are ways people naturally compare models. However, now they are all in one place and can be done systematically. These five ways are suggestions and have significant potential for expansion, especially as more models are added to \greensloth{}. Many ideas embedded in the ecosystem are heavily inspired by the models already included. A good example of how a comparison can be expanded on is to change the simulation examples from static, pre-calculated examples to a more interactive, dynamic, server-side simulation. This would allow the user not only to compare the models' results but also to test their own parameters and ideas. This idea is already being worked on, separately from \greensloth{}, but shows the potential to be added to the ecosystem~\cite{ComputationalBiologyAachenMxlweb2026}.

Overall, the \greensloth{} website is designed to be easy to understand and accessible to everyone, while also showcasing the models in their entirety. Because the documentation for each model is done in the same format, it is very clear when information is missing. This fact can be used to force the authors of models to provide all the necessary information, while helping them in the process through the structures of the \greensloth{} ecosystem. The website is not meant to replace a publication, but instead support it, by providing a dedicated platform for the models, where they are held to a specific standard, not in a sense of their value to the community, but in a sense of transparency, availability, and reusability. In turn, helping the model and its publication to follow the \gls{fair} and \gls{cure}. The website is a key way to support the sharing of models of photosynthesis, which may otherwise be lost in the onslaught of publications.

\subsection{Conclusion}

The \greensloth{} ecosystem is a solution for a wide variety of problems in the field of photosynthesis modelling. It provides support in creating model documentation to make the model's details as transparent as possible. It also provides ideas for demonstrations to make comparisons between models easier and reduce the user's effort. Finally, it provides an accessible, easy-to-use website to share and display models in a compact format. The website is also set up to compare each model side by side, providing another level of help for choosing models. 

Five models of photosynthesis have been chosen to show the power of \greensloth{}. They first had to be reimplemented, which was validated by the recreation of the figures in the publications. Then the complete documentation was completed to the \greensloth{} standard, including a summary, the model components, recreated figures, and demonstrations. These demonstrations are common but unofficial ways to compare models, methods that modellers already use to present their models. These demonstrations are standardized to allow at least a basic level of comparison between models and to reduce the effort required or even inspire users of the models. The models are implemented and shared openly on GitHub in a reproducible and clear way, along with the code for the demonstrations and figure recreation, and their documentation. Finally, all the information about these models is also shared openly on the \greensloth{} website, including additional ways to compare the models, such as sample simulations over time. This allows the user to access these five models of photosynthesis more interactively without detracting from the publication itself. The \greensloth{} ecosystem must not be seen as a replacement for the publication system or as a way to bypass it; it is merely a way to help model authors curate and share their work with the community.

The reimplementation of the five models was not entirely successful, as documentation is often misrepresented or missing. However, the overall trend is that the better the documentation, the more successful the reimplementation is. It does not matter how complex the model is; if documentation is taken care of, it can be reimplemented. This shows the importance of the \greensloth{} ecosystem, which provides a clear, easy-to-follow guide for documenting a model, with tools and features that support the author. The proposed demonstrations keep comparisons at a very basic level; however, this is necessary, as even with a focus on kinetic models of photosynthesis, the scope of models can vary dramatically. Keeping the demonstrations basic and more proof-of-concept-like allows them to encapsulate as many models as possible and show key differences between them. Still, the demonstrations shown at the current state need to be revised and worked on to allow even the basics of modelling to understand and use them. At the moment, some do not show the extreme cases well, such as the Li2021 model in the \gls{mca} demonstration~\figref{fig:mca-demon}, and others are too dependent on the author's expertise in the system, such as the fitting demonstration. These demonstrations are not perfect, but as more models are added to the \greensloth{} ecosystem, they will be expanded and improved. Additionally, more will be added as new inspiration is found from other models.

The website has been built with ease of access and use in mind for a wide range of users. The design is simple and modern, allowing users already familiar with the internet to feel comfortable with the website. However, the website has been designed only for use on a computer, without considering that a large part of internet users are on mobile devices~\cite{Digital2025Global2025}. This is a clear issue, as it limits the website's accessibility and, by extension, the models'. Additionally, as not many models are included on the website, the list is manageable; therefore, the search function has not been expanded. However, as more and more models are added, the search for models can quickly become as confusing as it is with publications; therefore, a more advanced search function is needed. The tag system is a good start; however, at its current state, it is too minimal and lacks sufficient information to be useful. More categories need to be considered, for example, how the models were simulated or which data were used to validate them. The side-by-side comparison page allows the user to compare two models directly on a specific aspect; however, it can be expanded further.

The \greensloth{} ecosystem is a project that incorporates many branches, all connected by a central idea. Making the creation, presentation, and sharing of models of photosynthesis easier and more accessible. The current state of the ecosystem presented in this thesis is only a starting block, like a sapling. This project has significant potential to grow and expand, with more models added, more demonstrations created, and the website improved. As a tree needs a crow to remove pests, fungi to refresh the soil, or a squirrel to plant its seeds, so does \greensloth{} need the help of the community to grow and thrive. A web developer to continue adapting the website, a web designer to make the website more appealing, a modeller to add more models, and a software developer to help with code and documentation. This additional help will not only let \greensloth{} grow, but will cement the community around it. How the ecosystem is built, which branches grow, which fall off, the main bark of \greensloth{} will keep being:

\begin{center}
    \Large
    \greensloth[{\textit{Photosynthesis Models at your Pace}}]
\end{center}